\documentclass[11pt,letterpaper]{article}

% ── Fonts ─────────────────────────────────────────────────────────
\usepackage[T1]{fontenc}
\usepackage{newtxtext}
\usepackage{amssymb}
\usepackage{newtxmath}
\usepackage[scaled=0.85]{inconsolata}

% ── Layout ────────────────────────────────────────────────────────
\usepackage[margin=1in]{geometry}
\usepackage{setspace}
\setstretch{1.08}
\usepackage{parskip}
\setlength{\parskip}{0.5\baselineskip plus 2pt minus 1pt}

% ── Section headings ──────────────────────────────────────────────
\usepackage{titlesec}
\titleformat{\section}{\large\bfseries}{\thesection}{1em}{}
\titleformat{\subsection}{\normalsize\bfseries}{\thesubsection}{0.8em}{}
\titleformat{\subsubsection}{\normalsize\itshape}{\thesubsubsection}{0.8em}{}
\titlespacing*{\section}{0pt}{1.8ex plus 0.5ex}{0.8ex plus 0.2ex}
\titlespacing*{\subsection}{0pt}{1.4ex plus 0.3ex}{0.4ex plus 0.1ex}

% ── Math ──────────────────────────────────────────────────────────
\usepackage{amsmath}
\usepackage{mathtools}
% newtxmath already defines \openbox; amsthm redefines it and errors out. Releasing the name before
% loading amsthm is the documented fix and changes nothing else.
\let\openbox\relax
\usepackage{amsthm}

\theoremstyle{definition}
\newtheorem{claim}{Claim}[section]
\newtheorem{exercise}{Exercise}[section]

% ── Tables ────────────────────────────────────────────────────────
\usepackage{booktabs}
\usepackage{array}
\usepackage{tabularx}
\renewcommand{\arraystretch}{1.15}
\usepackage{caption}
\captionsetup{font=small, labelfont=bf, labelsep=period, skip=6pt}
\captionsetup[table]{position=above}

% ── Floats ───────────────────────────────────────────────────────
\usepackage{float}
\usepackage{graphicx}
\usepackage{placeins}
\setlength{\floatsep}{6pt plus 2pt minus 2pt}
\setlength{\textfloatsep}{8pt plus 2pt minus 2pt}
\setlength{\intextsep}{6pt plus 2pt minus 2pt}
\renewcommand{\topfraction}{0.9}
\renewcommand{\bottomfraction}{0.9}
\renewcommand{\textfraction}{0.07}
\renewcommand{\floatpagefraction}{0.7}

% ── Lists ─────────────────────────────────────────────────────────
\usepackage{enumitem}
\setlist{nosep, leftmargin=1.5em}

% ── Colors & Links ────────────────────────────────────────────────
\usepackage{xcolor}
\definecolor{linkblue}{HTML}{0B5394}
\definecolor{citegreen}{HTML}{1B7340}
\definecolor{urlgray}{HTML}{444444}
\definecolor{boxgray}{HTML}{F4F4F2}
\definecolor{rulegray}{HTML}{BFBFBF}

% ── Boxes for the "why this matters" asides ──────────────────────
\usepackage[framemethod=TikZ]{mdframed}
\newmdenv[
  backgroundcolor=boxgray,
  linecolor=rulegray,
  linewidth=0.4pt,
  innertopmargin=6pt, innerbottommargin=6pt,
  innerleftmargin=8pt, innerrightmargin=8pt,
  skipabove=8pt, skipbelow=8pt,
  roundcorner=2pt,
]{aside}

% ── TikZ (diagrams) ──────────────────────────────────────────────
\usepackage{tikz}
\usetikzlibrary{arrows.meta, positioning, shapes.geometric, calc, decorations.pathreplacing, fit}

\usepackage[
  colorlinks=true,
  linkcolor=linkblue,
  citecolor=citegreen,
  urlcolor=urlgray,
  bookmarksnumbered=true,
  pdfstartview=FitH,
]{hyperref}
\usepackage[noabbrev,capitalise]{cleveref}

% cleveref does not auto-detect amsthm environments reliably, and the failure is SILENT: the label
% resolves (the number is right) but the name slot renders "??", and LaTeX logs no warning because
% nothing is actually undefined. Grepping the log is therefore useless here — `make pdf` greps the
% rendered PDF instead.
\crefname{claim}{Claim}{Claims}
\Crefname{claim}{Claim}{Claims}
\crefname{exercise}{Exercise}{Exercises}
\Crefname{exercise}{Exercise}{Exercises}

\usepackage{microtype}

% ── Notation macros ──────────────────────────────────────────────
\newcommand{\R}{\mathbb{R}}
\newcommand{\E}{\mathbb{E}}
\newcommand{\Var}{\operatorname{Var}}
\newcommand{\Cov}{\operatorname{Cov}}
\newcommand{\KL}{\mathrm{KL}}
\newcommand{\pol}{\pi_\theta}
\newcommand{\polref}{\pi_{\mathrm{ref}}}
\newcommand{\grad}{\nabla_\theta}
\newcommand{\score}{\grad \log \pol(y \mid x)}
\newcommand{\scorei}{g_i}
\newcommand{\Adv}{A}
\newcommand{\nsr}{\mathrm{NSR}}
\DeclareMathOperator*{\argmin}{arg\,min}

\title{\textbf{From REINFORCE to GRPO\\[4pt]
\large A worked tutorial, with the four things that actually bit us}}
\author{\texttt{assay} --- Phase 0.1 companion document}
\date{2026-08-02}

\begin{document}
\maketitle

\begin{abstract}\noindent
This document builds GRPO from REINFORCE one modification at a time, deriving each step rather than
asserting it, and stops at every point where the derivation makes a promise that our own measurements
broke. It is written to be read linearly. \Cref{sec:setup,sec:reinforce,sec:baseline,sec:group} are
the ladder: the objective, the score-function estimator, the baseline, and the group baseline that
gives GRPO its name. \Cref{sec:norm,sec:kl} add the two remaining pieces of full GRPO.
\Cref{sec:ladder,sec:ablations} map all of it onto the seven-run ladder and four ablations of Phase
0.1. \Cref{sec:surprises} is the part worth the most: four places where the standard account is
either incomplete or wrong at our scale, each one found by a measurement in this phase rather than
taken from the literature.

Every numerical claim in \cref{sec:ablations,sec:surprises} comes from the clean ladder run of
2026-08-02 --- ten arms, one commit, one GPU tier, 200 steps each --- and regenerates from committed
artefacts.
\end{abstract}

\tableofcontents

% ══════════════════════════════════════════════════════════════════
\section{Setup: what problem are we actually solving?}
\label{sec:setup}

\subsection{The objects}

Fix a pretrained language model with parameters $\theta$. Write:

\begin{itemize}
  \item $x$ --- a \emph{prompt}, drawn from a task distribution $\mathcal{D}$. In Phase 0.1, $x$ is a
        three-digit addition problem such as \texttt{"What is 471 + 638?"}.
  \item $y = (y_1, \dots, y_T)$ --- a \emph{completion}, a sequence of $T$ tokens.
  \item $\pol(y \mid x) = \prod_{t=1}^{T} \pol(y_t \mid x, y_{<t})$ --- the policy. This is just the
        language model's own next-token distribution, applied autoregressively.
  \item $R(x, y) \in \R$ --- the \emph{reward}, computed by a grader after the completion is
        finished. In Phase 0.1, $R \in \{0, 1\}$: one if the final integer in $y$ equals the correct
        answer, zero otherwise.
\end{itemize}

The objective is the expected reward:
\begin{equation}
  J(\theta) \;=\; \E_{x \sim \mathcal{D}} \; \E_{y \sim \pol(\cdot \mid x)} \big[ R(x, y) \big].
  \label{eq:objective}
\end{equation}

We want $\grad J(\theta)$ so we can do gradient \emph{ascent}.

\subsection{Why this is a bandit, not an MDP}

Textbook RL sets up states, actions, transitions and discounted returns. Almost none of that
machinery is needed here, and carrying it around obscures what is going on.

The reason: \textbf{reward arrives once, at the end, and the ``environment'' is deterministic.} The
model emits a completion; a grader looks at it; a number comes back. There is no stochastic
transition to model, no intermediate reward to discount, no value function that has to bootstrap
across timesteps.

Formally this is a \emph{contextual bandit}: context $x$, action $y$ (a whole completion, chosen from
an astronomically large action space), scalar payoff $R(x,y)$. Everything below is the bandit policy
gradient. When you see $\Adv_i$ called an ``advantage,'' remember it is not the $Q - V$ of temporal
difference learning; it is just a centred reward.

\begin{aside}
\textbf{Why this matters for reading the code.} \texttt{loop.py} assigns one scalar advantage per
\emph{rollout}, shared by all its tokens, and there is no per-timestep credit assignment anywhere.
That is not a simplification --- it is what the bandit formulation implies. Terminal reward and
deterministic transitions mean every token in a completion is equally ``responsible'' for the
outcome, as far as this estimator is concerned.
\end{aside}

\subsection{Vocabulary: rollout, episode, group, batch}

These get used loosely in the literature. In this project they mean:

\begin{itemize}
  \item \textbf{rollout} (= \textbf{episode}): one sampled completion $y$ for one prompt $x$, with
        its reward. These are synonyms.
  \item \textbf{group}: the $G$ rollouts sampled from the \emph{same} prompt $x$. Phase 0.1 uses
        $G = 8$. The group is the unit that makes GRPO work, so it is never split.
  \item \textbf{batch}: all rollouts used for one gradient step. Phase 0.1 uses $16$ prompts
        $\times\ 8$ rollouts $= 128$ rollouts per step.
\end{itemize}

% ══════════════════════════════════════════════════════════════════
\section{REINFORCE: the score-function estimator}
\label{sec:reinforce}

\subsection{The derivation}

We want $\grad J$. The difficulty is that $\theta$ appears inside the distribution we are averaging
over, not inside the thing being averaged --- $R$ does not depend on $\theta$ at all. So we cannot
simply push the gradient inside the expectation.

Fix $x$ and write $J_x(\theta) = \sum_y \pol(y \mid x) R(x,y)$. Then
\begin{align}
  \grad J_x(\theta)
    &= \sum_y \grad \pol(y \mid x) \, R(x,y) \label{eq:pg1}\\
    &= \sum_y \pol(y \mid x) \, \frac{\grad \pol(y \mid x)}{\pol(y \mid x)} \, R(x,y) \label{eq:pg2}\\
    &= \sum_y \pol(y \mid x) \, \grad \log \pol(y \mid x) \, R(x,y) \label{eq:pg3}\\
    &= \E_{y \sim \pol}\big[ R(x,y) \, \grad \log \pol(y \mid x) \big]. \label{eq:pg4}
\end{align}

Step \eqref{eq:pg1} $\to$ \eqref{eq:pg2} multiplies and divides by $\pol$, which is legal wherever
$\pol > 0$. Step \eqref{eq:pg2} $\to$ \eqref{eq:pg3} uses $\grad \log f = \grad f / f$ --- the
\emph{log-derivative trick}, and the entire content of the derivation. It converts a gradient of a
probability into an expectation, which is the one thing we can estimate by sampling.

Taking the outer expectation over prompts:
\begin{equation}
  \boxed{\;\grad J(\theta) \;=\; \E_{x, y}\big[ R(x,y) \, \score \big].\;}
  \label{eq:pgt}
\end{equation}

The quantity $\score$ is called the \emph{score function}. \Cref{eq:pgt} is the policy gradient
theorem for the bandit case.

\subsection{The estimator}

\Cref{eq:pgt} is an expectation, so replace it with a sample mean. Draw $N$ rollouts
$(x_i, y_i)$, write $r_i = R(x_i, y_i)$ and $\scorei = \grad \log \pol(y_i \mid x_i)$. Then
\begin{equation}
  \hat{g}_{\text{REINFORCE}} \;=\; \frac{1}{N}\sum_{i=1}^{N} r_i \, \scorei
  \label{eq:reinforce}
\end{equation}
is an unbiased estimator of $\grad J$. Ascend it. That is REINFORCE, complete.

In practice we do not compute $\scorei$ and then multiply --- we build a scalar surrogate loss whose
gradient is \eqref{eq:reinforce} and hand it to autograd:
\begin{equation}
  \mathcal{L} \;=\; -\frac{1}{N}\sum_{i=1}^{N} r_i \, \log \pol(y_i \mid x_i),
  \qquad \grad \mathcal{L} = -\hat{g}.
  \label{eq:surrogate}
\end{equation}
The minus sign turns ascent into descent. Note that $r_i$ is a \emph{constant} here: it is detached,
never differentiated through. If you differentiate through the reward you are doing something else
entirely.

\subsection{The identity that everything later depends on}

\begin{claim}[The score has mean zero]
\label{claim:zeromean}
For any $x$, $\;\E_{y \sim \pol(\cdot\mid x)}\big[\score\big] = 0.$
\end{claim}

\begin{proof}
\begin{equation}
  \E_y\big[\score\big]
  = \sum_y \pol(y\mid x) \grad \log \pol(y \mid x)
  = \sum_y \grad \pol(y \mid x)
  = \grad \sum_y \pol(y \mid x)
  = \grad 1 = 0. \qedhere
\end{equation}
\end{proof}

This is not a technicality. It is the load-bearing fact underneath baselines
(\cref{sec:baseline}), and \cref{sec:surprises} is largely the story of what happens when an
innocuous implementation detail breaks it.

\begin{aside}
\textbf{Read \cref{claim:zeromean} intuitively.} $\score$ points in the direction that would make
\emph{this particular completion} more likely. Averaged over completions actually drawn from the
policy, those directions cancel exactly --- because the policy is already normalised, and making
some outputs more likely necessarily makes others less so. There is no free ``push everything up''
direction. Keep this picture; \cref{sec:lennorm} is about an implementation that quietly creates one.
\end{aside}

% ══════════════════════════════════════════════════════════════════
\section{The variance problem, and baselines}
\label{sec:baseline}

\subsection{Why REINFORCE is unusable as written}

$\hat{g}_{\text{REINFORCE}}$ is unbiased but has enormous variance. The cleanest way to see the
problem: with $R \in \{0,1\}$, \cref{eq:reinforce} is
\begin{equation}
  \hat{g} = \frac{1}{N}\sum_{i \,:\, r_i = 1} \scorei,
\end{equation}
a sum over the \emph{correct} rollouts only. Failures contribute exactly nothing. Two consequences:

\begin{enumerate}
  \item \textbf{Half the data is discarded.} At a $47\%$ pass rate, $53\%$ of the compute produces no
        gradient signal.
  \item \textbf{Every surviving term has the same sign.} The estimator can only ever say ``do more of
        this,'' never ``do less of that.'' It is not contrastive.
\end{enumerate}

\subsection{The baseline, and why it is free}

Subtract a constant $b$ from every reward:
\begin{equation}
  \hat{g}_b \;=\; \frac{1}{N}\sum_{i=1}^{N} (r_i - b)\, \scorei.
  \label{eq:baseline}
\end{equation}

\begin{claim}[Baselines are unbiased]
\label{claim:unbiased}
If $b$ does not depend on $y$, then $\E[\hat{g}_b] = \grad J$ for every $b$.
\end{claim}

\begin{proof}
$\E[(R - b)\score] = \E[R \score] - b\,\E[\score] = \grad J - b \cdot 0 = \grad J$, by
\cref{claim:zeromean}.
\end{proof}

So $b$ is a free parameter: it changes the variance of the estimator without changing what the
estimator estimates. This is the single most useful fact in policy-gradient methods.

The effect on the update is immediate and worth stating plainly. With $b > 0$, a failure
($r_i = 0$) now carries advantage $-b$, so its term is $-b \cdot \scorei$ --- pushing the policy
\emph{away} from that completion. \textbf{The estimator becomes contrastive.} Both consequences of
\cref{sec:baseline} are addressed by one subtraction.

\subsection{How much variance does it remove?}
\label{sec:varcalc}

Take $R \in \{0,1\}$ with $\Pr[R = 1] = p$, and define the conditional second moments of the score
\begin{equation}
  E_c = \E\big[\|\score\|^2 \,\big|\, R = 1\big], \qquad
  E_w = \E\big[\|\score\|^2 \,\big|\, R = 0\big].
\end{equation}
Since all estimators here share the same mean (\cref{claim:unbiased}), comparing variances is the
same as comparing second moments. For constant $b$,
\begin{equation}
  \E\big[(R-b)^2 \|\score\|^2\big] \;=\; p(1-b)^2 E_c \;+\; (1-p)\,b^2 E_w.
  \label{eq:secondmoment}
\end{equation}

Minimising over $b$ gives the optimal baseline
\begin{equation}
  b^\star \;=\; \frac{p\,E_c}{p\,E_c + (1-p)E_w}.
  \label{eq:optb}
\end{equation}

\begin{aside}
\textbf{\Cref{eq:optb} is not $\bar{r}$.} The optimal baseline equals the mean reward $p$ only in the
special case $E_c = E_w$ --- when correct and incorrect completions carry equally large score
vectors. Everyone uses $\bar{r}$ anyway, because it is trivial to compute. \Cref{sec:eweight} is
about what our own measurements say when that assumption fails, and it turns out to fail badly.
\end{aside}

Setting $E_c = E_w$ and comparing $b = 0$ against $b = p$:
\begin{align}
  b = 0: &\quad p\,E, \\
  b = p: &\quad p(1-p)^2 E + (1-p)p^2 E \;=\; E\,p(1-p)\big[(1-p) + p\big] \;=\; E\,p(1-p),
\end{align}
so the variance ratio is
\begin{equation}
  \boxed{\;\frac{\Var[\hat{g}_0]}{\Var[\hat{g}_p]} \;=\; \frac{p E}{E p (1-p)} \;=\; \frac{1}{1-p}.\;}
  \label{eq:ratio}
\end{equation}

At $p = 0.5$ the baseline halves the variance; at $p = 0.9$ it cuts it by a factor of ten. This is
the pre-registered point prediction for ablation A (\cref{sec:abA}).

% ══════════════════════════════════════════════════════════════════
\section{From a global baseline to a group baseline}
\label{sec:group}

\subsection{The idea GRPO is named after}

\Cref{claim:unbiased} allows $b$ to depend on the \emph{prompt} $x$, since the expectation in
\cref{claim:zeromean} is taken at fixed $x$. So we may use $b(x)$, and the variance-minimising choice
is approximately
\begin{equation}
  b(x) \;\approx\; \E\big[R \mid x\big],
\end{equation}
the expected reward \emph{on this prompt}. This is strictly better than a global $b$: an easy prompt
and a hard prompt have very different expected rewards, and a single global constant is systematically
wrong for both.

Classical actor--critic methods estimate $\E[R \mid x]$ with a learned value network $V_\phi(x)$ ---
a second model, with its own parameters, its own optimiser, and its own failure modes.

\textbf{GRPO's contribution is to replace that network with a Monte Carlo estimate.} Sample $G$
completions for the same prompt and average their rewards:
\begin{equation}
  \hat{b}(x) \;=\; \frac{1}{G}\sum_{j=1}^{G} r_j.
  \label{eq:groupmean}
\end{equation}

That is the whole idea. \emph{Group Relative Policy Optimisation} is REINFORCE with a per-prompt
baseline bought with extra samples instead of a critic network.

\begin{aside}
\textbf{The trade being made.} A critic amortises: train it once, query it cheaply forever. A group
baseline pays $G\times$ the sampling cost at every step, and buys in return: no second model, no
critic-training instability, no value-function bias, and an estimate that is automatically correct
for the current policy. At $G=8$ you spend $8\times$ the rollouts to delete an entire neural network
from the system. For LLM post-training, where sampling is the well-understood part and critic
training is not, that is usually a good trade.
\end{aside}

\subsection{Leave-one-out}

\Cref{eq:groupmean} has a subtle flaw: $r_i$ appears in its own baseline, so
$\hat{b}$ is correlated with $r_i$ and \cref{claim:unbiased} no longer applies exactly. The fix is to
exclude the sample being centred:
\begin{equation}
  \Adv_i^{\text{loo}} \;=\; r_i - \frac{1}{G-1}\sum_{j \neq i} r_j.
  \label{eq:loo}
\end{equation}

\begin{claim}[LOO is a uniform rescale of the mean-centred advantage]
\label{claim:loo}
With $S = \sum_j r_j$ and $\Adv_i^{\text{mean}} = r_i - S/G$,
\begin{equation}
  \Adv_i^{\text{loo}} \;=\; \frac{G\,r_i - S}{G-1} \;=\; \frac{G}{G-1}\,\Adv_i^{\text{mean}}.
\end{equation}
\end{claim}

\begin{proof}
$\Adv_i^{\text{loo}} = r_i - \frac{S - r_i}{G-1} = \frac{(G-1)r_i - S + r_i}{G-1} = \frac{G r_i - S}{G-1}$,
and $\Adv_i^{\text{mean}} = \frac{G r_i - S}{G}$.
\end{proof}

At $G = 8$ the factor is $8/7 \approx 1.143$. So LOO differs from mean-centring by a constant scale
--- which is absorbed entirely by the learning rate, and \emph{exactly} cancelled if you normalise by
the standard deviation (\cref{sec:norm}). The bias argument for LOO is real but small; do not expect
it to show up as a visible effect.

\subsection{Worked example: what conditioning actually buys}
\label{sec:worked}

This is the table that makes the difference concrete. Take $G = 8$, a batch whose overall pass rate is
$\bar{r} = 0.43$, and three groups of differing difficulty.

\begin{table}[H]
\centering
\caption{The same rewards, three ways of centring them. Only the \emph{group} baseline reacts to how
hard the prompt is.}
\label{tab:worked}
\small
\begin{tabular}{@{}llr@{}}
\toprule
\textbf{Group} & \textbf{Baseline} & \textbf{Advantages $\Adv_i$} \\
\midrule
Easy: $r = (1,1,1,1,1,1,1,0)$
  & \texttt{none}      & $+1,\,+1,\,+1,\,+1,\,+1,\,+1,\,+1,\;\;0$ \\
  & \texttt{global}    & $+0.57 \times 7,\;\; -0.43$ \\
  & \texttt{group\_loo}& $+0.143 \times 7,\;\; \mathbf{-1.000}$ \\
\midrule
Hard: $r = (0,0,0,0,0,0,0,1)$
  & \texttt{none}      & $0 \times 7,\;\; +1$ \\
  & \texttt{global}    & $-0.43 \times 7,\;\; +0.57$ \\
  & \texttt{group\_loo}& $-0.143 \times 7,\;\; \mathbf{+1.000}$ \\
\midrule
Average: $r = (1,0,1,1,0,0,1,0)$
  & \texttt{none}      & $+1,\;0,\;+1,\;+1,\;0,\;0,\;+1,\;0$ \\
  & \texttt{global}    & $+0.570 / -0.430$ \\
  & \texttt{group\_loo}& $+0.571 / -0.571$ \\
\bottomrule
\end{tabular}
\end{table}

Read the three blocks in order:

\begin{itemize}
  \item \textbf{Easy prompt.} Under \texttt{none}, seven rollouts on a problem the model
        \emph{already solves} each receive full weight $+1$: the update is dominated by reinforcing
        existing competence. Under \texttt{group\_loo} they get $+0.143$, while the single failure
        gets $-1.000$ --- the strongest signal in the group. The group baseline has correctly
        identified that the only informative event here is the failure.
  \item \textbf{Hard prompt.} Under \texttt{none} the seven failures receive advantage exactly
        $0$: a failure carries no information at all. Under \texttt{group\_loo} the lone success gets
        the maximum $+1.000$.
  \item \textbf{Average prompt.} \texttt{global} and \texttt{group\_loo} nearly coincide
        ($0.570$ vs $0.571$) --- because this group's pass rate equals the batch's.
\end{itemize}

\begin{aside}
\textbf{The one-sentence summary of \cref{sec:group}.} The group baseline differs from the global one
\emph{only} when a prompt is harder or easier than average --- and that difference \emph{is} the
conditioning. If every prompt in your task set had identical difficulty, rung 2 and rung 3 would be
the same algorithm.
\end{aside}

% ══════════════════════════════════════════════════════════════════
\section{Advantage normalisation}
\label{sec:norm}

Published GRPO divides the centred advantage by the group's standard deviation:
\begin{equation}
  \tilde{\Adv}_i \;=\; \frac{\Adv_i}{\sigma_{\text{group}} + \varepsilon},
  \qquad \sigma^2_{\text{group}} = \frac{1}{G}\sum_j \big(r_j - \bar{r}\big)^2 .
  \label{eq:znorm}
\end{equation}
Combined with mean-centring, this makes $\tilde{\Adv}$ a \emph{$z$-score} within the group.

Two properties follow, and both matter operationally.

\begin{claim}[Bounded]
\label{claim:bound}
For any group of $G$ values, $\;\max_i |\tilde{\Adv}_i| \le \sqrt{G-1}$.
\end{claim}

\begin{proof}
The $z$-scores satisfy $\sum_i \tilde{\Adv}_i = 0$ and $\sum_i \tilde{\Adv}_i^2 = G$. Fix $i=1$ and
write $z = \tilde{\Adv}_1$. Then $\sum_{i \ge 2} \tilde{\Adv}_i = -z$ and
$\sum_{i\ge2}\tilde{\Adv}_i^2 = G - z^2$. Cauchy--Schwarz gives
$\big(\sum_{i\ge2}\tilde{\Adv}_i\big)^2 \le (G-1)\sum_{i\ge2}\tilde{\Adv}_i^2$, i.e.
$z^2 \le (G-1)(G - z^2)$, hence $z^2 G \le G(G-1)$ and $|z| \le \sqrt{G-1}$.
\end{proof}

At $G = 8$ this is $\sqrt{7} \approx 2.6458$, attained exactly by a lone winner (or lone loser).
Phase 0.1 logs \texttt{max\_abs\_advantage} every step purely as a rig check: if it ever exceeds
$2.6458$, the advantage function is not computing a $z$-score and nothing downstream means anything.

\begin{claim}[Scale-invariant]
\label{claim:scaleinv}
$\tilde{\Adv}$ is unchanged by any affine transform $r_i \mapsto \alpha r_i + \beta$, $\alpha \ne 0$.
\end{claim}

This one has a sharp consequence that we got wrong on the first pass and had to correct before
running anything. \Cref{claim:scaleinv} says the rewards $\{0,\dots,0,1\}$ and
$\{0.5,\dots,0.5,0.500001\}$ produce \emph{identical} normalised advantages. Normalisation keeps only
the \emph{shape} of reward differences and discards their \emph{magnitude} entirely. We return to
this in \cref{sec:abC}, because it is the whole content of ablation C.

\subsection{Dr.\ GRPO's objection}

Two conventions are contested in the literature. Published GRPO (i) divides by $\sigma_{\text{group}}$
and (ii) length-normalises the log-probability sum. Dr.\ GRPO argues both introduce bias and removes
them. Phase 0.1 keeps both as switches --- \texttt{normalize\_by\_std} and \texttt{length\_normalize}
--- precisely so the question is answered by measurement rather than by citation. \Cref{sec:lennorm}
reports what we found, and it is not a small effect.

% ══════════════════════════════════════════════════════════════════
\section{The KL leash}
\label{sec:kl}

Nothing so far stops the policy from drifting arbitrarily far from where it started in pursuit of
reward. The standard remedy adds a penalty against a frozen reference policy $\polref$ (the model as
it was before RL):
\begin{equation}
  \mathcal{L} \;=\; \underbrace{-\frac{1}{N}\sum_i \Adv_i \log \pol(y_i \mid x_i)}_{\text{policy gradient}}
  \;+\; \beta \underbrace{\KL\big[\pol \,\|\, \polref\big]}_{\text{leash}}.
  \label{eq:klloss}
\end{equation}

\subsection{Estimating the KL}

We cannot compute $\KL$ in closed form over the space of completions, so we estimate it from the same
samples. The estimator used here is Schulman's \emph{k3}:
\begin{equation}
  \hat{k}_3 \;=\; \rho - \log \rho - 1, \qquad \rho = \frac{\polref(y \mid x)}{\pol(y \mid x)} .
  \label{eq:k3}
\end{equation}
It is unbiased --- $\E_{\pol}[\rho] = 1$ and $\E_{\pol}[\log \rho] = -\KL[\pol \| \polref]$, so
$\E[\hat{k}_3] = 1 - (-\KL) - 1 = \KL$ --- and, unlike the naive $-\log\rho$, it is non-negative
pointwise, which makes it far better behaved as a penalty.

\begin{aside}
\textbf{Apply k3 per token, not per sequence.} Sequence-level probabilities are products over
hundreds of tokens, so $\rho$ compounds multiplicatively and its variance explodes with length.
Phase 0.1 computes \eqref{eq:k3} at each token position and averages. This was a deliberate change
during implementation, not an accident of the code.
\end{aside}

\subsection{Two implementation notes}

\textbf{Add it to the loss, not the reward.} If the KL penalty is folded into $R$ before advantages
are computed, it passes through the group baseline --- and a penalty that is roughly constant within
a group is \emph{subtracted away} by that baseline. The leash silently goes slack. It must enter at
\eqref{eq:klloss}, after centring.

\textbf{Under LoRA, the reference is free.} With a frozen base model and only LoRA adapters trainable,
$\polref$ is exactly the base model --- obtained by disabling the adapter. No second set of weights
is needed, only a second forward pass.

% ══════════════════════════════════════════════════════════════════
\section{The ladder}
\label{sec:ladder}

Everything above is now one dataclass. Each rung differs from the previous one by exactly one field,
which is what makes any observed difference attributable.

\begin{table}[H]
\centering
\caption{The Phase 0.1 ladder. Runs 4--6 are cut: under the pinned single-epoch design the importance
ratio is identically $1$, so clipping is a no-op, and the remaining two switches are isolated by
ablations B and C instead.}
\label{tab:ladder}
\small
\begin{tabular}{@{}llll@{}}
\toprule
\textbf{Run} & \textbf{Name} & \textbf{Configuration} & \textbf{What it adds} \\
\midrule
1 & REINFORCE           & \texttt{baseline="none"}       & --- \\
2 & $+$ mean baseline   & \texttt{baseline="global"}     & contrastive updates \\
3 & $+$ group baseline  & \texttt{baseline="group\_loo"} & conditioning on the prompt \\
7 & full GRPO           & $+$ \texttt{normalize\_by\_std}, \texttt{kl\_coef} & scale-free advantages, a leash \\
\bottomrule
\end{tabular}
\end{table}

The rung $1 \to 2 \to 3$ progression is worth stating carefully, because it is easy to conflate the
last two steps:

\begin{itemize}
  \item $1 \to 2$ introduces \textbf{centring}. Failures acquire negative weight; the update becomes
        contrastive.
  \item $2 \to 3$ introduces \textbf{conditioning}. The centre becomes prompt-specific, so an easy
        prompt and a hard prompt are judged against their own difficulty rather than the batch's
        average. \Cref{tab:worked} is exactly this difference.
\end{itemize}

% ══════════════════════════════════════════════════════════════════
\section{The four ablations}
\label{sec:ablations}

Phase 0.1's deliverable is four \emph{deliberate breakages}, each with a pre-registered signature
recorded before any training run, and each carrying two distinct failure branches: \emph{prediction
wrong} (a finding, to be reported) and \emph{rig broken} (a bug, to be fixed). All numbers below come
from the clean ladder of 2026-08-02, one commit, one GPU tier, seed 0.

\subsection{Ablation A --- remove the baseline}
\label{sec:abA}

\textbf{Isolates:} the baseline itself, rung 1 versus rung 2.
\textbf{Prediction:} \cref{eq:ratio} --- the no-baseline estimator should have $1/(1-p)$ times the
noise-to-signal ratio of the baselined one.

This ablation is the subject of \cref{sec:cosine,sec:lennorm,sec:eweight}, because both the original
measurement and the theory behind it turned out to be wrong. It is the most instructive of the four.

\subsection{Ablation B --- remove the KL leash on a degenerate reward}
\label{sec:abB}

\textbf{Isolates:} $\beta$. Because B needs \emph{two} changes from run 7 (no leash \emph{and} a
degenerate reward), it is measured against a control that has the degenerate reward but keeps the
leash; the difference between B and its control isolates $\beta$ alone.

The degenerate reward $R_{\text{format}}$ scores a completion on emitting the right \emph{shape},
regardless of whether the answer is correct. It is a deliberately hackable grader.

\begin{table}[H]
\centering
\caption{Ablation B. Proxy is what is optimised; true is $R_{\text{binary}}$, measured and never
trained on. The gap is the Goodhart quantity.}
\label{tab:abB}
\small
\begin{tabular}{@{}lrrrrr@{}}
\toprule
\textbf{Arm} & \textbf{proxy} & \textbf{true} & \textbf{gap} & \textbf{KL end} & \textbf{KL share of loss} \\
\midrule
control ($\beta = 0.04$) & 0.993 & 0.486 & $+0.507$ & 2.68 & 0.490 \\
ablation B ($\beta = 0$) & 0.999 & 0.510 & $+0.489$ & 2.45 & 0.000 \\
\bottomrule
\end{tabular}
\end{table}

Two things to take from \cref{tab:abB}. First, the degenerate grader is \emph{comprehensively} hacked:
the policy scores $0.993$ on the proxy while true performance sits at $0.486$ --- it learned the
format and not the task. That $+0.507$ gap is the largest in the whole ladder. Second, and against the
pre-registered prediction, \textbf{removing the leash changes the gap by $-0.018$, i.e.\ nothing} ---
despite the KL term carrying $49\%$ of the loss in the control arm.

The mechanism turns out to run through ablation D: the degenerate reward drives the dead-group
fraction to $0.888$, so there is almost no policy gradient left for the leash to restrain. A leash is
irrelevant when the horse has stopped.

\subsection{Ablation C --- add a tiny tie-breaker}
\label{sec:abC}

\textbf{Isolates:} the reward's magnitude structure. Replace $R_{\text{binary}}$ with
\begin{equation}
  R_{\text{tiebreak}}(x,y) \;=\; R_{\text{binary}}(x,y) \;+\; 0.001 \cdot |y| .
\end{equation}

The pre-registered signature for this ablation was originally ``divide-by-$\approx 0$ advantage
spikes.'' \Cref{claim:bound,claim:scaleinv} say that is \textbf{unreachable}: the normalised advantage
is a $z$-score, bounded by $\sqrt{G-1}$ and invariant to scale. We rewrote the signature before
running anything.

What actually goes wrong is sharper, and is this project's thesis in miniature. Because $z$-scoring
discards magnitude, a group that was unanimous under the binary reward --- and therefore contributed
\emph{zero} gradient --- becomes, on adding a $0.001$ tie-breaker, \textbf{fully alive at full
gradient magnitude}, with the advantages ordered entirely by completion length.

\begin{table}[H]
\centering
\caption{Ablation C against its parent. Goodhart emerging from the optimiser's own arithmetic.}
\label{tab:abC}
\small
\begin{tabular}{@{}lrrrr@{}}
\toprule
\textbf{Arm} & \textbf{true reward} & \textbf{gap} & \textbf{dead groups} & \textbf{tokens/rollout} \\
\midrule
run 7 (binary)     & $0.592 \to 0.907$ & $+0.000$ & 0.468 & 17.2 \\
ablation C (tiebreak) & $0.598 \to 0.728$ & $+0.042$ & \textbf{0.008} & \textbf{35.5} \\
\bottomrule
\end{tabular}
\end{table}

All four predicted signatures fire: dead groups collapse $0.468 \to 0.008$, completions double, the
true-reward gain falls from $+0.315$ to $+0.130$, and a real proxy--true gap of $+0.042$ opens on a
task where run 7's gap is identically zero.

\subsection{Ablation D --- force every group unanimous}
\label{sec:abD}

\textbf{Isolates:} the dead-group failure mode, by construction.

\begin{claim}[Dead-group probability]
\label{claim:dead}
If a prompt has per-sample success probability $p$, the probability that a group of $G$ i.i.d.\
rollouts is unanimous --- and therefore produces \emph{zero} gradient under a group baseline --- is
\begin{equation}
  P_{\text{dead}}(p) \;=\; p^G + (1-p)^G .
\end{equation}
\end{claim}

This function is convex, symmetric about $p = 1/2$, and minimised there. At $G = 8$ it is $0.008$ at
$p = 0.5$ but $0.66$ at $p = 0.05$ \emph{or} $p = 0.95$. Two consequences that are easy to miss:

\begin{itemize}
  \item \textbf{It is blind to sign.} A group that is unanimously correct (saturated) and one that is
        unanimously wrong (unreachable) are both dead, and behave oppositely under training --- one
        will stay dead, one may become live. The scalar cannot tell you which.
  \item \textbf{The mean hides it.} By Jensen's inequality, $\E[P_{\text{dead}}(p)] \ge
        P_{\text{dead}}(\E[p])$ for convex $P_{\text{dead}}$. A task set that is half trivial and half
        impossible has mean pass rate $0.5$ and \emph{every} group dead; a task set genuinely centred
        at $0.5$ has $0.8\%$ dead. Identical means, a $55\times$ difference in wasted compute. This is
        why Phase 0.1's task was selected on the dead-group fraction rather than on pass rate.
\end{itemize}

Measured: with every group forced unanimous, \texttt{frac\_degenerate\_groups} $= 1.000$ on all $200$
steps, \texttt{grad\_norm} $= 0.0000$ on all $200$ steps, true reward flat at $0.503 \to 0.509$, and
$313{,}312$ tokens generated at the parent's rate. \textbf{Thirty-eight minutes of GPU for exactly
zero gradient.}

\subsection{C and D are one event with two faces}

The two ablations are the same phenomenon, separated by whether the reward has a tie-breaking term:

\begin{table}[H]
\centering
\caption{What a unanimous group does, as a function of reward structure.}
\label{tab:cd}
\small
\begin{tabular}{@{}ll@{}}
\toprule
\textbf{Unanimous group, reward is\ldots} & \textbf{Result} \\
\midrule
clean binary              & zero gradient --- the step is wasted (D) \\
carrying any tie-breaker  & full-magnitude gradient chasing noise (C) \\
\bottomrule
\end{tabular}
\end{table}

% ══════════════════════════════════════════════════════════════════
\section{Four things the standard account gets wrong at this scale}
\label{sec:surprises}

This section is the reason the document exists. Each subsection is a place where the clean derivation
above makes a promise that our own measurements broke.

\subsection{Dead groups are a pathology GRPO \emph{introduces}}
\label{sec:deadintro}

A group is dead when all its advantages are zero. Under which baselines does that happen?

\begin{itemize}
  \item \texttt{none}: $\Adv_i = r_i$, so the group is dead only if \emph{every} rollout failed.
  \item \texttt{global}: $\Adv_i = r_i - \bar{r}_{\text{batch}}$. Since $\bar{r}$ is a batch mean and
        $r_i$ is binary, $\Adv_i$ is essentially \emph{never} exactly zero. Dead groups cannot occur.
  \item \texttt{group\_loo}: $\Adv_i = 0$ for all $i$ exactly when the group is unanimous, either way.
\end{itemize}

Measured on the clean ladder:

\begin{table}[H]
\centering
\caption{Dead-group fraction by baseline, averaged over 200 steps. The analytical prediction, exactly.}
\label{tab:dead}
\small
\begin{tabular}{@{}llr@{}}
\toprule
\textbf{Baseline} & \textbf{Dead when\ldots} & \textbf{Measured} \\
\midrule
\texttt{none}       & the group is all-wrong           & 0.037 \\
\texttt{global}     & never                            & \textbf{0.000} \\
\texttt{group\_loo} & the group is unanimous, either way & \textbf{0.454} \\
\texttt{forced} (D) & always, by construction          & 1.000 \\
\bottomrule
\end{tabular}
\end{table}

\textbf{Read the middle two rows against each other.} The same unanimous group that GRPO throws away
is perfectly alive under a global baseline, where it contributes $\pm b$. Conditioning the baseline on
the prompt --- the thing that makes GRPO better --- is \emph{also} what creates the dead-group
pathology.

And it \emph{worsens as training succeeds}: under \texttt{group\_loo} the dead fraction climbs from
$0.244$ over the first ten steps to $0.606$ over the last ten, because the policy is solving more
prompts unanimously. Full GRPO (run 7) is worse still, $0.300 \to 0.694$. The better the policy gets,
the larger the share of each batch that produces no gradient at all --- a built-in brake that tightens
exactly as the task is being learned.

This is not a curiosity: \texttt{prime-rl}, Prime Intellect's production RL stack, ships a
\texttt{zero\_advantage} pre-batch filter \emph{on by default} to drop exactly these rollouts. The
pathology is recognised industry-wide. What we add is measuring its size per configuration.

% Short form for the ToC: \cref does not survive a moving argument and renders as "??" there.
\subsection[Length normalisation breaks the zero-mean score]
           {Length normalisation breaks \cref{claim:zeromean}}
\label{sec:lennorm}

Implementations divide each rollout's summed log-probability by its token count, so that a 40-token
rollout does not contribute four times the gradient of a 10-token one:
\begin{equation}
  \mathcal{L} = -\frac{1}{N}\sum_i \Adv_i \cdot \frac{1}{|y_i|}\log \pol(y_i \mid x_i).
  \label{eq:lennorm}
\end{equation}

This looks innocuous. It is not. \Cref{claim:zeromean} relied on the sum telescoping,
$\sum_y \grad\pol = \grad\sum_y \pol = \grad 1 = 0$. With a per-sample weight $w(y) = 1/|y|$ inside
the sum, that telescoping fails:
\begin{equation}
  \E_y\Big[\tfrac{1}{|y|}\score\Big] \;=\; \sum_y \pol(y\mid x)\,\tfrac{1}{|y|}\,\grad\log\pol(y\mid x)
  \;=\; \sum_y \tfrac{1}{|y|}\grad \pol(y \mid x) \;\neq\; \grad\!\!\sum_y \pol(y\mid x) \;=\; 0,
\end{equation}
because $1/|y|$ cannot be pulled out of the sum. \textbf{The score no longer has mean zero}, so
\cref{claim:unbiased} fails and a baseline is no longer free: it changes the estimand.

We measured this directly with a paired fixed-policy probe --- one policy state, 40 batches, the
\emph{same} rollouts scored under all three baselines, so nothing but the baseline differs. The
statistic is the cosine between the three arms' mean gradients: if all three estimate the same thing,
these should be $\approx 1$.

\begin{table}[H]
\centering
\caption{Do the three baselines estimate the same gradient? Mean over 3 seeds, $N=40$ batches.}
\label{tab:lennorm}
\small
\begin{tabular}{@{}lrrr@{}}
\toprule
 & \texttt{none}\,$|$\,\texttt{global} & \texttt{none}\,$|$\,\texttt{group\_loo} & \texttt{global}\,$|$\,\texttt{group\_loo} \\
\midrule
\texttt{length\_normalize=True}  & $+0.727$ & $+0.734$ & $+0.996$ \\
\texttt{length\_normalize=False} & $\mathbf{+0.991}$ & $\mathbf{+0.992}$ & $+0.999$ \\
\bottomrule
\end{tabular}
\end{table}

With length normalisation on, the two \emph{centred} estimators agree with each other almost exactly
($0.996$) while the uncentred one sits roughly $40^\circ$ away from both. Turn length normalisation
off and all three snap into agreement. \Cref{claim:zeromean} is restored, and with it the premise that
makes ablation A a well-posed question at all.

\begin{aside}
\textbf{This has teeth beyond ablation A.} The ladder's primary arms all run with
\texttt{length\_normalize=True}. In that configuration the rungs are not estimating a common quantity,
so ``does the baseline reduce variance'' has no well-defined answer --- \emph{that} is the finding,
not a failure to measure. Separately, removing length normalisation improved the noise-to-signal ratio
of \emph{every} arm ($0.45$--$0.94 \to 0.30$--$0.33$) and improved final true reward on both arms
tested ($0.907 \to 0.927$ and $0.728 \to 0.844$), with the proxy--true gap unchanged. Two independent
lines of support for Dr.\ GRPO's position, arrived at by measurement.
\end{aside}

\subsection{The variance reduction does not materialise}
\label{sec:eweight}

With the estimators made comparable (\texttt{length\_normalize=False}), we can finally ask ablation
A's actual question. Define the scale-free noise-to-signal ratio for baseline $b$,
\begin{equation}
  \nsr_b \;=\; \frac{V_b}{\|\bar{g}_b\|^2},
  \qquad V_b = \frac{1}{N}\sum_{n=1}^{N} \big\| g_n^{(b)} - \bar{g}_b \big\|^2 ,
\end{equation}
where $g_n^{(b)}$ is the full-batch gradient of batch $n$ under baseline $b$. Scale-freeness is
required, not cosmetic: the arms produce gradients of very different magnitudes, so comparing raw
variance would report that scale difference instead.

\Cref{eq:ratio} predicts $\nsr_{\texttt{none}} / \nsr_{\texttt{global}} = 1/(1-p)$. At the measured
$p = 0.466$ that is $1.87$.

\begin{table}[H]
\centering
\caption{Observed ratio against the point prediction, $N=40$ batches, paired bootstrap.}
\label{tab:nsr}
\small
\begin{tabular}{@{}rrrr@{}}
\toprule
\textbf{Seed} & \textbf{ratio} & \textbf{95\% CI} & \textbf{predicted $1/(1-p)$} \\
\midrule
0 & 0.859 & $[0.616,\ 1.219]$ & 1.869 \\
1 & 0.921 & $[0.704,\ 1.155]$ & 1.825 \\
2 & 1.350 & $[0.969,\ 1.873]$ & 1.922 \\
\bottomrule
\end{tabular}
\end{table}

Every interval contains $1.0$ --- no reduction is demonstrated --- and every interval \emph{excludes}
the prediction. Mean ratio $1.04$. \textbf{At $p \approx 0.47$ on this task, the baseline delivers no
detectable variance reduction, and the textbook magnitude is ruled out.}

Why? Recall from \cref{sec:varcalc} that \cref{eq:ratio} assumed $E_c = E_w$. Solve
\cref{eq:secondmoment} for the ratio at which the benefit vanishes entirely:
\begin{align}
  p E_c &= p(1-p)^2 E_c + (1-p)p^2 E_w \nonumber\\
  E_c\big[1 - (1-p)^2\big] &= (1-p)\,p\,E_w \nonumber\\
  E_c \, p(2-p) &= (1-p)\,p\,E_w \nonumber\\
  \boxed{\;\frac{E_w}{E_c} = \frac{2-p}{1-p} \;=\; 2.87 \ \text{ at } p = 0.466.\;}
  \label{eq:ewec}
\end{align}

So if incorrect completions carry roughly $2.9\times$ the squared score norm of correct ones, the
baseline's benefit disappears exactly. That is entirely plausible --- an incorrect completion is one
the policy assigned lower probability to, and $\|\grad\log\pol\|$ is larger there. \Cref{eq:ewec} is a
\emph{derived, falsifiable number}, and measuring $E_w/E_c$ directly is the obvious next probe.

\begin{aside}
\textbf{The methodological point, which outlives the result.} On the first pass, with length
normalisation on, the observed ratio was $1.786$ against a prediction of $1.872$ --- a $4.6\%$ match,
and we nearly reported it as a clean confirmation. Decomposing $\nsr$ showed the agreement was
coincidental: theory predicts \emph{signal unchanged, noise down $1.87\times$}, whereas what actually
happened was \emph{noise up $2.34\times$, signal up $4.16\times$}. Two effects the theory does not
contain, pulling in opposite directions, landing on the right number. The rig check --- ``do these
three estimators even share a mean?'' --- is the only thing that caught it.
\end{aside}

\subsection{The half-batch cosine cannot compare baseline arms}
\label{sec:cosine}

Ablation A was originally measured by splitting each batch in half, computing both half-gradients, and
taking their cosine. Two independent samples of the same expected gradient should agree; the cosine
is then a direct read on estimator noise, with no trend confound.

It came back \emph{reversed}: $\rho_2/\rho_1 = 0.154$ against a pre-registered gate of $\ge 2.0$, with
all rig checks clean. The reason is that both arms are confounded, in opposite directions, and neither
confound is about variance.

\begin{itemize}
  \item \texttt{baseline="none"} --- every advantage is $\ge 0$, so \emph{both} halves weight the
        shared component of $\score$ positively. The cosine is \textbf{inflated} by precisely the
        nuisance direction a baseline exists to remove.
  \item \texttt{baseline="global"} --- the baseline is the \emph{full-batch} mean while the cosine
        splits that same batch. Writing $b = (b_A + b_B)/2$ for the two half-means, each half carries
        a term $\pm (b_A - b_B)/2 \cdot \sum \score$: an anti-correlated component
        \textbf{manufactured by the split boundary}. The cosine is \textbf{deflated}.
  \item \texttt{baseline="group\_*"} --- advantages sum to zero \emph{within} each group and groups
        are never split, so each half sums to zero independently. No coupling. This arm was always
        clean.
\end{itemize}

The obvious repair --- centre each half on its own mean advantage before measuring --- does not work,
and the reason is worth internalising. Under \texttt{none} the centred advantage is
$r_i - \bar{r}_A$; under \texttt{global} it is $(r_i - b) - (\bar{r}_A - b) = r_i - \bar{r}_A$. The
constant cancels: \textbf{centring makes the two arms the same estimator}, erasing the contrast
instead of the confound. Verified numerically --- at a shared policy state the two agree to six
decimal places.

That is why ablation A needs the fixed-policy probe of \cref{sec:eweight} rather than a comparison
between training arms. Two training arms follow different trajectories, so any difference between them
confounds the estimator with the policy state it was measured at. Only a probe that holds the policy
fixed and varies \emph{nothing but the baseline} can answer the question.

\begin{aside}
\textbf{The general lesson.} The half-batch cosine measures directional \emph{reproducibility of the
applied update}. A degenerate estimator that always points the same way scores near $1.0$ while
learning nothing from reward. So a metric that rewards a large common-mode component will always
report that removing it hurt --- which is Goodhart's law operating inside the diagnostic itself. For a
project whose entire subject is the gap between a target and the proxy used to measure it, finding
that gap in our own instrument is not an embarrassment. It is the thesis.
\end{aside}

% ══════════════════════════════════════════════════════════════════
\section{Summary}
\label{sec:summary}

\begin{enumerate}
  \item \textbf{REINFORCE} is the log-derivative trick plus a sample mean: $\hat{g} = \frac{1}{N}\sum
        r_i \scorei$. Unbiased, unusably noisy, and non-contrastive when $R \ge 0$.
  \item \textbf{A baseline} is free because $\E[\score] = 0$. It makes the update contrastive and, in
        theory, cuts variance by $1/(1-p)$.
  \item \textbf{GRPO's group baseline} is a Monte Carlo estimate of $\E[R \mid x]$ bought with $G$
        rollouts instead of a critic network. It differs from a global baseline exactly when prompts
        differ in difficulty.
  \item \textbf{Normalisation} makes advantages $z$-scores: bounded by $\sqrt{G-1}$ and
        scale-invariant. The scale-invariance is what makes ablation C possible.
  \item \textbf{The KL leash} uses the per-token k3 estimator and must be added to the loss, not the
        reward, or the group baseline subtracts it away.
\end{enumerate}

And the four things measurement changed:

\begin{enumerate}
  \item Dead groups are a pathology the \emph{group} baseline introduces --- $0.000$ under a global
        baseline, $0.454$ under GRPO's.
  \item Length normalisation breaks $\E[\score] = 0$, so under the ladder's primary configuration the
        rungs do not estimate a common quantity at all.
  \item Once they do, the predicted variance reduction \textbf{does not appear} --- consistent with
        $E_w/E_c \approx 2.9$, a derived and testable number.
  \item The half-batch cosine cannot compare baseline arms, and the obvious fix erases the contrast
        rather than the confound.
\end{enumerate}

\end{document}
